\section{Corrections to the statements of \texorpdfstring{\cite{Baek}}{Baek's paper} that are used here}\label{app:corrections}

The audit that accompanies the formalization compared every result of~\cite{Baek} with its TeX source and with the formal statement~\cite{Companion}. Every result of the paper
holds with the intended statement. \cref{tab:corrections} lists the corrections that concern the results used in this paper, with the labels E$n$ of the audit. The complete list of the audit has
27 items, E22 only pointing to E21, and two missing hypotheses: \baek{Thm.~2.1.1} (Schneider's Theorem~4.2.3: $\sigma_K(X)$ is the length of $\bigcup_{t\in X}e_K(t)$) needs
$K$ to have interior points, and \baek{Thm.~3.1.2} needs a bounded Nef polygon; every use satisfies them (the text itself uses $\sigma_K$ only as defined in \cref{fact:convex}\ref{fact:convex:sigma}, which does not depend on Thm.~2.1.1). Twelve of the findings in the audit come from the notes of the formalization~\cite{Cureton}.

{\small
\begin{longtable}{@{}>{\raggedright\arraybackslash}p{0.065\linewidth}>{\raggedright\arraybackslash}p{0.19\linewidth}>{\raggedright\arraybackslash}p{0.37\linewidth}>{\raggedright\arraybackslash}p{0.22\linewidth}@{}}
\caption{Corrections to~\cite{Baek} that concern the results used here.}\label{tab:corrections}\\
\toprule
Item & Result of~\cite{Baek} & Correction & Where here\\
\midrule\endfirsthead
\toprule
Item & Result of~\cite{Baek} & Correction (continued) & Where here\\
\midrule\endhead
\bottomrule\endlastfoot
E1 & Def.~2.3.9 & the left and right side of a line are defined for lines not parallel to the $x$-axis (printed: $y$-axis) & \cref{fact:monotone}\ref{fact:monotone:envelope}\\
E2 & Thm.~2.4.1, 2.4.2 & the union over $t\in[0,\omega]$ and over $(0,\omega)$ agree, since $F_\omega\cap Q^-_K(0)=F_\omega\cap Q^-_K(\omega)=\emptyset$; $\cC(S)$ is bounded also for $\omega=\pi/2$ & \cref{fact:monotone}\ref{fact:monotone:cap}\\
E3 & Prop.~2.5.4 & the reflection exchanges the walls ($a\leftrightarrow c$, $b\leftrightarrow d$) and the ends $W\leftrightarrow Z$; $K$ and not $K^{\mir}$ stands on the right of the vertex identity & \cref{fact:monotone}\ref{fact:monotone:mirror}, \cref{lem:triangle}, \cref{prop:curvature}\\
E4 & Lemma~2.5.6, Thm.~2.5.8 (and Lemma~3.4.8, Remark~3.4.1) & used without proof, true for every cap: $A_K^-(0)=(h_K(0),0)$, $C_K^+(\omega)=h_K(\omega+\pi/2)v_\omega$, and $O\in K$ if $\omega<\pi/2$; the same holds for $o_\omega$, which Lemma~3.4.8 and Remark~3.4.1 use (not in the audit) & \cref{lem:ends}, \cref{lem:contains}, \cref{lem:wedge}, \cref{lem:sandwich}\\
E5 & Thm.~3.1.2, Lemma~3.4.7 & the Nef polygon must be bounded; the two applications of the theorem for a pinned normal change disjoint strips & \cref{fact:push}\\
E6 & Def.~3.2.5 & $\cN_\Theta(K)=F_\omega\cap\bigcup_{t\in\Theta}Q^-_K(t)$ (printed: $P_\omega$ in place of $F_\omega$); with $P_\omega$, Prop.~3.3.5 and Lemma~3.4.2 are false & \cref{sec:polygon}, \cref{fact:width}, \cref{prop:penmax}\\
E7 & Prop.~3.3.1\,(2) & normal angles in $\Theta^\diamond\cup\{\omega+\pi,3\pi/2\}$ (printed: $\Theta^\diamond$) & \cref{sec:polygon}\\
E8 & Thm.~3.4.3, 3.5.2, 3.5.4 & minor gaps: a Hausdorff limit of polygon caps is a polygon cap (resp.\ a cap), and the niche of the limit lies in the polygon niches eventually, point by point & \cref{prop:penmax}, \cref{prop:selection}, \cref{lem:limsup}\\
E9 & Lemma~4.2.4 & the range is $\omega\in[\sec^{-1}(2.2),\pi/2)$ (printed: $[\tan^{-1}(2.2),\pi/2)$) & \cref{fact:corner}\\
E10 & Thm.~4.2.5, 1.5.2 & the width bound for the pentagon $P_\omega\setminus\Delta_\omega$ is asserted without proof; the proof puts the vertices of $\Delta_\omega$ in the open quarter-plane & \cref{lem:triangle}, \cref{lem:width}\\
E12 & Thm.~6.1.2 & the proof (Gerver's Thm.~2 makes $\Gs$ a balanced maximum sofa) is invalid; the statement is true and follows from Romik's equations & \cref{fact:gerver}\ref{fact:gerver:cap}\\
E13 & Prop.~6.2.2 & $f^\pm_{K^{\mir}}(t)=g^\mp_K(\pi/2-t)$, and likewise for $g$ (printed: without the change of angle) & \cref{fact:arms}\ref{fact:arms:mirror}\\
E14 & Def.~6.2.2, Thm.~6.2.3 & $\partial^+$ is the right derivative (Def.~6.2.2 calls it the left one), and the second display of Thm.~6.2.3 is about $\partial^-\xx_K$ & \cref{fact:arms}\ref{fact:arms:deriv}\\
E15 & Lemma~6.3.1 & cites Thm.~3.5.4 (balanced maximum caps) for $\cN(K)\subset K$ of a maximum polygon cap; the applicable result is Thm.~3.4.10, which gives what the lemma needs, and the lemma holds & not used here (the box of \cref{prop:penmax} bounds the diameter)\\
E17 & Lemma~6.5.2--6.5.5, Thm.~6.5.6 & $f_K(0)=1$ is used without being stated; Lemma~6.5.5 is used on $(0,\pi/2]$ (printed: $(0,1]$) & \cref{fact:arms}\ref{fact:arms:regular}, \cref{fact:iteration}\ref{fact:iteration:eleven}\\
E20 & Lemma~8.3.5 & the arc $\mathbf u_K^{0,\pi}$ is outside the range $b<a+\pi$ of \baek{Thm.~7.3.2}; the area is split at the atom $\pi/2$ instead & \cref{fact:Q}\ref{fact:Q:decomp}\\
E21 & Lemma~8.1.7, Thm.~8.2.4 & $\cN(K)\subset K$ is used but not available on $\Ki$; both statements hold and have proofs that avoid it & \cref{fact:Q}\ref{fact:Q:triple}\\
E23 & Lemma~8.3.6\,(3) & the tangent line parametrization is evaluated at $\pi/2$ (printed: at $\varphi^{\mathrm L}$, where the statement is false) & \cref{fact:Q}\ref{fact:Q:decomp}\\
E24 & Thm.~8.4.1 & stated without proof; true, and proved from Romik's equations by interval arithmetic; part (4) holds with one-sided derivatives at the junctions and the ends of the phases. Not in the audit: part (2), read as a Jordan region, omits the open floor segment of $\cN(\KG)$, and the exact description is the identity in \cref{fact:gerver}\ref{fact:gerver:niche} & \cref{fact:gerver}\\
E25, E26 & Thm.~8.4.3\,(2), Prop.~8.4.4\,(4) & $\bB(t_3)$ (printed: $\bD(t_3)$); the measure $\breve\sigma_D$ lives on $(\pi/2+t_0,\pi/2+t_2]$ with a shifted density & \cref{fact:Q}\ref{fact:Q:max}\\
E27 & several & misprints: the second half-plane of $Q^-_S(t)$ lacks its ``$-1$'' (Prop.~2.2.2); the sign of the change of $\abs{\cN_\Theta}$ in the proof of Lemma~3.4.7; $\cR_B$ starts at $\pi+\varphi^{\mathrm R}$, not $\pi/2+\varphi^{\mathrm R}$ (Def.~8.3.3); the statement of Thm.~2.1.3 & \cref{sec:sofas}, \cref{fact:convex}, \cref{fact:push}, \cref{sec:Q}\\
\end{longtable}
}
